\section{Conclusion}\label{sec:conclusion}

This note isolates a simple price consequence of directional access friction in a spatial duopoly. When the right-hand firm becomes more distant from a fixed rival, two forces move together: ordinary spatial separation softens competition, but the same displacement increases the directional burden borne by consumers who buy from that firm. On an open parameter region of a strict global price Nash equilibrium, the second force is strong enough to reverse the moving firm's price response. Its equilibrium price falls while the rival's price rises.

The result is deliberately a pricing result rather than a theory of endogenous location. It also does not claim robustness to arbitrary asymmetric transport functions. Its contribution is to show, in a transparent covered Hotelling environment, that one-sided directional exposure can overturn the usual own-price effect of greater separation without overturning the rival's response. Setting the directional term to zero restores the standard quadratic benchmark in which both prices rise with separation.